\documentclass[11pt]{article}
\usepackage[margin=1in]{geometry}
\usepackage{amsmath,amssymb,amsthm,mathtools}
\usepackage[colorlinks=true,linkcolor=blue,citecolor=blue,urlcolor=blue]{hyperref}

\newtheorem{theorem}{Theorem}
\newtheorem{proposition}{Proposition}
\newtheorem{remark}{Remark}
\newcommand{\D}{D_{\mathcal P_N}}
\newcommand{\norm}[2]{\left\lVert #1\right\rVert_{#2}}

\title{A supercritical exponential-Orlicz lower gain for discrepancy}
\author{Autonomous solution packet for arXiv:1411.5794}
\date{11 August 2026}

\begin{document}
\maketitle

\begin{abstract}
Bilyk and Markhasin ask for lower bounds matching their exponential-Orlicz
upper estimates for discrepancy in dimensions $d\geq3$.  The critical
$\exp(L^{2/(d-1)})$ lower bound is already Roth's $L_2$ bound.  We extract
additional information from the short Riesz-product proof of Bilyk--Lacey--
Vagharshakyan: their $L_1$ test function has enough $L_2$ control to give a
growing-$p$ lower bound.  Consequently, for every $\beta>2$, the
$\exp(L^\beta)$ discrepancy of every $N$-point set is larger than the Roth
order by a positive power of $\log N$ (up to one $\log\log N$ factor).  This
is a partial result; it does not reach the conjectural matching exponent.
\end{abstract}

\section{Statement}

For an $N$-point set $\mathcal P_N\subset[0,1)^d$, use the unnormalised
anchored discrepancy
\[
 D_{\mathcal P_N}(x)
 =\sum_{z\in\mathcal P_N}{\bf1}_{[0,x)}(z)-N x_1\cdots x_d.
\]
As usual, $\norm{f}{\exp(L^\beta)}$ denotes the Luxemburg norm associated to
a Young function equal to $e^{t^\beta}-1$ for large $t$.  On a probability
space,
\begin{equation}\label{eq:orlicz-moments}
 \norm{f}{\exp(L^\beta)}\asymp_\beta
 \sup_{p>1}p^{-1/\beta}\norm{f}{L^p}.
\end{equation}

\begin{theorem}[supercritical Orlicz gain]\label{thm:main}
For every $d\geq3$ there is $\eta_d>0$ such that, for every $\beta>2$,
there is $c_{d,\beta}>0$ with the following property.  Every $N$-point set
$\mathcal P_N\subset[0,1)^d$, $N\geq3$, satisfies
\begin{equation}\label{eq:main}
 \norm{D_{\mathcal P_N}}{\exp(L^\beta)}
 \geq c_{d,\beta}
 (\log N)^{\frac{d-1}{2}+\eta_d(1-2/\beta)}
 (\log\log(e^eN))^{-1/\beta}.
\end{equation}
One may take $\eta_d=\varepsilon_d/4$, where $\varepsilon_d>0$ is any
admissible short-product exponent in the proof of Bilyk--Lacey--
Vagharshakyan's improved star-discrepancy lower bound.
\end{theorem}

For every fixed $\beta>2$, the exponent in \eqref{eq:main} is strictly larger
than Roth's $(d-1)/2$.  As $\beta\to\infty$, it tends to the star-discrepancy
gain $(d-1)/2+\eta_d$ obtained by the same short Riesz product.

\section{The growing-moment extraction}

Fix $n$ by $2N\leq2^n<4N$.  We recall only the quantitative properties of
the test function constructed in Section 9 of
Bilyk--Lacey--Vagharshakyan (BLV).  Choose a sufficiently small
$\varepsilon=\varepsilon_d>0$, put
\[
 q=\lfloor a n^\varepsilon\rfloor,\qquad b=\frac14,
 \qquad \widetilde\rho=a q^b n^{-(d-1)/2},
\]
and use the separated blocks of multiindices in the discrepancy part of their
proof.  The blocks have length comparable to $n/q$; this is also the length
used explicitly in their estimates following the displayed block definition.
For the BLV $r$-functions $f_{\mathbf r}$ let
$F_t=\sum_{\mathbf r\in\mathbb A_t}f_{\mathbf r}$ and form the short product
\[
 \Psi=\prod_{t=1}^q(1+\widetilde\rho F_t)
      =1+\Psi^{\mathrm{sd}}+\Psi^{\neg}.
\]
Here $\Psi^{\mathrm{sd}}$ retains the products whose multiindices are strongly
distinct and $\Psi^{\neg}$ is the coincidence remainder.

The BLV proof gives
\begin{align}
 \norm{\Psi^{\mathrm{sd}}}{L^1}&\leq C_d,
 \label{eq:l1}\\
 \langle D_{\mathcal P_N},\Psi^{\mathrm{sd}}\rangle
 &\geq c_d q^{1/4}n^{(d-1)/2}.
 \label{eq:pairing}
\end{align}
Indeed, the linear term has the size on the right of \eqref{eq:pairing},
whereas all higher strongly-distinct terms pair with discrepancy by a summable
amount.  We next record the extra estimate that is not extracted in BLV.

\begin{proposition}[near-$L^1$ control of the test]\label{prop:test}
For $q$ sufficiently large,
\begin{equation}\label{eq:l2test}
 \norm{\Psi^{\mathrm{sd}}}{L^2}
 \leq \exp\!\bigl(C_d q^{1/2}\log(2q)\bigr).
\end{equation}
Consequently, if $p\geq C_d q^{1/2}\log(2q)$ and
$p'=p/(p-1)$, then
\begin{equation}\label{eq:dualnorm}
 \norm{\Psi^{\mathrm{sd}}}{L^{p'}}\leq C_d.
\end{equation}
\end{proposition}

\begin{proof}
The proof uses estimates already established inside the BLV short-product
argument.  Their Lemma 5.3, uniformly for partial products, gives
\begin{equation}\label{eq:partial2}
 \left\lVert\prod_{t\in V}(1+\widetilde\rho F_t)\right\rVert_2
 \leq \exp(C_d q^{2b})=\exp(C_d q^{1/2}).
\end{equation}
Their proof also records the crude fourth-moment estimate
\begin{equation}\label{eq:partial4}
 \left\lVert\prod_{t\in V}(1+\widetilde\rho F_t)\right\rVert_4
 \leq (C_dq)^q,
\end{equation}
again uniformly in $V\subset\{1,\ldots,q\}$.

Write the inclusion--exclusion expansion of the coincidence remainder in the
form
\[
 \Psi^{\neg}=\sum_G A_G B_G,
\]
where $G$ runs over the admissible coincidence graphs,
$A_G=\widetilde\rho^{|V(G)|}\operatorname{SumProd}(X(G))$, and $B_G$
is a partial product over the unused blocks.  The estimates in BLV that prove
$\norm{\Psi^{\neg}}{L^1}\lesssim1$ actually establish the stronger summability
statement
\begin{equation}\label{eq:graphsum}
 \sum_G\norm{A_G}{L^s}\leq C_d,
 \qquad s=q^{2b}=q^{1/2}.
\end{equation}
For $s\geq4$, choose $r$ by $1/2=1/s+1/r$.  Thus
$r=2s/(s-2)\in[2,4]$.  Interpolating \eqref{eq:partial2} and
\eqref{eq:partial4}, with interpolation parameter $4/s$, yields
\begin{equation}\label{eq:partialr}
 \sup_V\left\lVert\prod_{t\in V}(1+\widetilde\rho F_t)\right\rVert_r
 \leq \exp\!\bigl(C_d q^{1/2}\log(2q)\bigr).
\end{equation}
H\"older, \eqref{eq:graphsum}, and \eqref{eq:partialr} now imply
\[
 \norm{\Psi^{\neg}}{L^2}
 \leq \sum_G\norm{A_G}{L^s}\norm{B_G}{L^r}
 \leq \exp\!\bigl(C_d q^{1/2}\log(2q)\bigr).
\]
Since $\Psi^{\mathrm{sd}}=\Psi-1-\Psi^{\neg}$ and
\eqref{eq:partial2} also bounds $\norm{\Psi}{L^2}$, this proves
\eqref{eq:l2test}.

Finally, interpolate \eqref{eq:l1} and \eqref{eq:l2test}.  Since
\[
 \frac1{p'}=1-\frac1p=(1-\theta)\cdot1+\theta\cdot\frac12,
 \qquad \theta=\frac2p,
\]
we get
\[
 \norm{\Psi^{\mathrm{sd}}}{L^{p'}}
 \leq C_d\exp\!\left(\frac{C_dq^{1/2}\log(2q)}p\right).
\]
This is bounded when $p\geq C_dq^{1/2}\log(2q)$.
\end{proof}

\begin{proposition}[growing-$p$ discrepancy lower bound]\label{prop:lp}
With the notation above,
\begin{equation}\label{eq:lp}
 \norm{D_{\mathcal P_N}}{L^p}
 \geq c_dq^{1/4}n^{(d-1)/2}
 \quad\text{whenever}\quad
 p\geq C_dq^{1/2}\log(2q).
\end{equation}
\end{proposition}

\begin{proof}
Apply H\"older to \eqref{eq:pairing} and use
\eqref{eq:dualnorm}.
\end{proof}

\section{Proof of Theorem \ref{thm:main}}

Take $p$ comparable to $q^{1/2}\log(2q)$ in
\eqref{eq:orlicz-moments}.  Proposition \ref{prop:lp} gives
\begin{align*}
 \norm{D_{\mathcal P_N}}{\exp(L^\beta)}
 &\gtrsim_{d,\beta}
 p^{-1/\beta}q^{1/4}n^{(d-1)/2}\\
 &\gtrsim_{d,\beta}
 n^{(d-1)/2}
 q^{1/4-1/(2\beta)}(\log(2q))^{-1/\beta}.
\end{align*}
Since $q\asymp n^{\varepsilon_d}$ and
$\eta_d=\varepsilon_d/4$, this is
\[
 n^{(d-1)/2+\eta_d(1-2/\beta)}
 (\log(2n))^{-1/\beta}.
\]
Finally $n\asymp\log N$.  Enlarging the constant handles the finitely many
small values for which $q<4$, and proves \eqref{eq:main}.

\section{Scope and obstruction to a full result}

The Bilyk--Markhasin upper bound for $\beta\geq2/(d-1)$ has order
$(\log N)^{d-1-1/\beta}$.  Theorem \ref{thm:main} is not matching.
At the critical exponent $2/(d-1)$, the Roth lower bound already gives the
correct order of their critical upper theorem; the open difficulty concerns
the stronger supercritical scale.

The threshold $\beta>2$ is intrinsic to the present extraction.  If one keeps
a general short-product parameter $b\leq1/4$, positivity and the second-moment
argument permit a gain $q^b$, while controlling the coincidence remainder in
$L^2$ through the available graph estimate costs moments at scale
$q^{1-2b}$ (up to logarithms).  The inequality
$b>(1-2b)/\beta$ is easiest at the extremal admissible value $b=1/4$ and then
becomes exactly $\beta>2$.  Reaching $\beta\leq2$, much less the full matching
scale, requires a genuinely stronger coincidence estimate or a different test
function.  Such an improvement is tied to the unresolved higher-dimensional
small-ball/star-discrepancy problem.

\begin{thebibliography}{9}
\bibitem{BM}
D. Bilyk and L. Markhasin,
\emph{BMO and exponential Orlicz space estimates of the discrepancy function
in arbitrary dimension}, arXiv:1411.5794.

\bibitem{BLV}
D. Bilyk, M. T. Lacey, and A. Vagharshakyan,
\emph{On the Small Ball Inequality in All Dimensions}, arXiv:0705.4619.

\bibitem{DHMP}
J. Dick, A. Hinrichs, L. Markhasin, and F. Pillichshammer,
\emph{Discrepancy of second order digital sequences in function spaces with
dominating mixed smoothness}, arXiv:1604.08713.
\end{thebibliography}

\end{document}
